\chapter{Prezentacja warstwy użytkowej projektu}

Rozdział poświęcony jest wizualnej prezentacji zrealizowanego projektu. Zanim przejdziemy do omówienia i prezentacji samych widoków strony, warto pochylić się nad ogólnymi założeniami, które definiowały ostateczny kształt samej strony.
\section{Identyfikacja wizualna i stylistyka strony}
Projekt strony utrzymany został w minimalistycznej estetyce, gdzie sama strona jest witryną typu \textit{single page}, co oznacza, że nie zawiera ona podstron.
\section{Kolorystyka}
W projekcie wykorzystano minimalistyczną paletę barw, gdzie kolory pochodzą ze strony \textit{\url{https://coolors.co/}}. Całość zapewnia dobrą czytelność niezależnie od warunków oświetleniowych. Dzięki zastosowaniu pseudoklasy \texttt{:root} w pliku \texttt{style.css}, kolory mogą być szybko zmienione.
\begin{itemize}
    \item \textbf{Tryb jasny (Light Mode):} Dominujący kolor - biel (\texttt{\#F8F9FA}) jako tło, jak i kolor paska nawigacyjnego, głęboka czerń (\texttt{\#0B090A}) dla tekstów i elementów interaktywnych oraz odcień szarości (\texttt{\#DEE2E6}), pełniącą rolę separatorów, będący również kolorem stopki.
    \item \textbf{Tryb ciemny (Dark Mode):} Odwrócona paleta barw, gdzie tło stanowi czerń, a teksty i ikony przybierają odcień bielszy, co redukuje ostatecznie zmęczenie wzroku przy słabym oświetleniu.
\end{itemize}
\section{Typografia}
W projekcie zastosowano kilka krojów z biblioteki \textit{Google Fonts}, które budują tożsamość wizualną strony:
\begin{itemize}
    \item \textbf{Cinzel:} Krój dekoracyjny inspirowany antyczną rzymską epigrafiką. Został wykorzystany dla nagłówków sekcji oraz wyników kalkulatora. Całość nadaje elegancki, historyczny charakter.
    \item \textbf{Montserrat:} Nowoczesny, bezszeryfowy krój pisma o wysokiej czytelności. Wykorzystany jako główny font treściowy (body), w listach zadań oraz formularzach.
\end{itemize}
\section{Prezentacja witryny}
Prezentacja wyglądu strony, wraz z przykładowym działaniem funkcjonalności, przedstawiona będzie w poniższych podsekcjach.
\subsection{Pasek nawigacji}
Wygląd paska nawigacji przedstawiono na zdjęciu poniżej - \textbf{rys. \ref{rys 5.1}}
\begin{figure}[H]
    \centering
    \includegraphics[scale=0.4]{figures/5.1.eps}
    \caption{Panel dodawania nowego użytkownika}\label{rys 5.1}
\end{figure}
\subsection{Sekcja Hero}
Wygląd sekcji Hero z nagłówkiem dekoracyjnym przedstawiono na zdjęciu poniżej - \textbf{rys. \ref{rys 5.2}}
\begin{figure}[H]
    \centering
    \includegraphics[scale=0.4]{figures/5.2.eps}
    \caption{Główna sekcja powitalna witryny (Hero)}\label{rys 5.2}
\end{figure}
\subsection{Lista To-Do}
Wygląd listy To-Do przedstawiono na zdjęciu poniżej - \textbf{rys. \ref{rys 5.3}}
\begin{figure}[H]
    \centering
    \includegraphics[scale=0.4]{figures/5.3.eps}
    \caption{Lista To-Do ze stworzonym zadaniem}\label{rys 5.3}
\end{figure}
\subsection{Kalkulator rzymski}
Kalkulator do konwersji liczb arabskich na liczby rzymskie przedstawiono na zdjęciu poniżej - \textbf{rys. \ref{rys 5.4}}
\begin{figure}[H]
    \centering
    \includegraphics[scale=0.4]{figures/5.4.eps}
    \caption{Kalkulator liczb rzymskich - 2026}\label{rys 5.4}
\end{figure}
\subsection{Formularz kontaktowy}
Formularz kontaktowy przedstawiony jest na zdjęciu poniżej - \textbf{rys. \ref{rys 5.5}}
\begin{figure}[H]
    \centering
    \includegraphics[scale=0.4]{figures/5.5.eps}
    \caption{Formularz kontaktowy}\label{rys 5.5}
\end{figure}
\subsection{Alert - Błąd}
Powiadomienie (górna część strony) informujące o błędzie w danych formularza przedstawione jest na rysunku poniżej - \textbf{rys. \ref{rys 5.6}}
\begin{figure}[H]
    \centering
    \includegraphics[scale=0.4]{figures/5.6.eps}
    \caption{Powiadomienie o błędzie}\label{rys 5.6}
\end{figure}
\subsection{Alert - Sukces}
Powiadomienie (górna część strony) informujące o sukcesie przesłania formularza przedstawione jest na rysunku poniżej - \textbf{rys. \ref{rys 5.7}}
\begin{figure}[H]
    \centering
    \includegraphics[scale=0.4]{figures/5.7.eps}
    \caption{Powiadomienie o powodzeniu przesłania formularza}\label{rys 5.7}
\end{figure}
\subsection{Tryb ciemny}
Strona internetowa w ciemnym motywie przedstawiona jest na zdjęciu poniżej - \textbf{rys. \ref{rys 5.8}}
\begin{figure}[H]
    \centering
    \includegraphics[scale=0.4]{figures/5.8.eps}
    \caption{Tryb ciemny strony internetowej}\label{rys 5.8}
\end{figure}